\documentclass[11pt,a4paper]{article}
\usepackage[margin=25mm]{geometry}
\usepackage[T1]{fontenc}
\usepackage{lmodern,amsmath,amssymb,amsthm,booktabs,xcolor}
\usepackage{xurl,microtype}
\usepackage[colorlinks=true,linkcolor=blue!45!black,urlcolor=blue!45!black]{hyperref}
\newtheorem{theorem}{Theorem}[section]
\newtheorem{lemma}[theorem]{Lemma}
\newtheorem{proposition}[theorem]{Proposition}
\newtheorem{corollary}[theorem]{Corollary}
\newcommand{\pos}[1]{\left(#1\right)_+}
\newcommand{\ind}{\mathbf 1}
\newcommand{\LP}{L_{\mathrm{raw}}}
\title{Erd\H{o}s 708: proof ladders for the unrestricted asymptotic\\
\large Executed steps, adaptive carriers, and the remaining arithmetic inequality}
\author{Research checkpoint}
\date{23 November 2026}
\begin{document}
\maketitle
\begin{abstract}
We pursue three routes to $g(n)\le(2+o(1))n$: adaptive Hall transport,
joint diffuse transfer, and a signed or cluster-based carrier certificate.
New proved steps include uniform one-sided concentration in every arithmetic
bank, a demand-scale transfer for one jointly clipped target budget, a
positive-density small-price reduction, and structural reductions of the
adaptive Hall cuts. Two actual arithmetic examples refute tempting stronger
intermediate assertions. A relaxation of the allowable complement from load
$1$ to load $1+\delta/2$ uses the asymptotic slack explicitly.
\textbf{The unrestricted asymptotic has not been proved.} The exact $2n$
conjecture is also not claimed. The last missing inequality, its parameter
order, and the reason each attempted shortcut fails are stated explicitly.
No new values of $g(n)$ are computed.
\end{abstract}

\section{Scope, retrieval, and the common endpoint}

The newest accessible dot input was the package quoted in the conversation;
the additional history and recent-file checks did not expose a newer dot
proof. This report continues that package and the previous audited checkpoint.
It does not claim access to a separate live dot session. Complete working
proofs accompany this report in the three route directories.

Fix an $M$-term positive interval $I$ and $c=2+\delta$, $0<\delta\le1$.
For $0\le w_p\le1$ put
\[
 X(a)=\sum_{p\le M}w_pv_p(a),\quad
 H_X(b)=\pos{X(b)-1},\quad
 J_c(w)=\sum_{a\le M}\pos{X(a)-c}-\sum_{b\in I}H_X(b).
\]
At a positive global maximum take the entire strict tail
$A=\{a\le M:X(a)>c\}$, $n=|A|$. The preceding checkpoint establishes
\begin{align}
 \max_wJ_c(w)&=\max_{S\subseteq[2,M]}(\LP(S,I)-c|S|),\label{global}\\
 F=\LP(A,I)&=cn+J_c(w),\quad z_b=1\quad(X(b)>1),\label{paired}\\
 \sum_{b\in I}z_bv_p(b)&=\sum_{a\in A}v_p(a)\quad(w_p>0),\label{rows}\\
 \sum_{a\in A}X(a)&\le C_\delta n,\quad
 \sum_p\frac{w_p}{p-1}\le B_\delta,\label{bounds}\\
 \sum_{a\in A}\exp\left(\frac{X(a)-c}{2B_\delta}\right)&\le2n.
 \label{mgf}
\end{align}
Increase $B_\delta$ to at least one if necessary. These are inherited inputs;
their notes are included with the new work. The earlier near-quadratic
argument excludes unbounded positive maxima with $M\ge\kappa n^2$ for any
fixed $\kappa>0$. The trivial interval cover and Wang's long-interval gate
leave $cn<M<16n^2$; the genuine unresolved branch is $M=o(n^2)$.
Wang's existing fractional-to-integer theorem is used only at the endpoint
\cite{WangRaw,WangMain}; rounding is not reproved here.

The exact next result need not have zero exceptional sources.
Let $1\le\tau<c-1$ and $\gamma=c-1-\tau>0$. Suppose a single target budget
pays fractions $\theta_a\in[0,1]$ of the demands $X(a)-(1+\tau)$, with a total
extra allowance $E$. Put $t=\sum_A(1-\theta_a)$. Weighted Jensen applied to
\eqref{mgf} gives
\begin{equation}
 \boxed{\quad J_c(w)\le
 2B_\delta t\log(2n/t)-\gamma(n-t)+E.\quad}              \label{finish}
\end{equation}
Indeed, the unserved excess above $c$ is at most
$2B_\delta t\log(2n/t)$; every served fraction contributes the gain $\gamma$.
For example, it suffices to have
\[
 \frac tn\le\frac{\gamma}{32B_\delta\log(64B_\delta/\gamma)},
 \qquad E\le\frac{\gamma n}{4}.
\]
The right side of \eqref{finish} is then at most $-3\gamma n/8$.
Thus a sufficiently small \emph{fixed} loss fraction, chosen after $\delta$,
is enough. Requiring an $o(n)$ exceptional set at every stage can be avoided.

If a uniform construction excludes all positive maxima with sufficiently
large $n$, the remaining positive values are bounded by a constant depending
only on $\delta$, using \eqref{bounds}. Equation \eqref{global} then yields
$\LP(S,I)\le(2+\delta)|S|+C_\delta'$. Letting $\delta$ tend to zero and applying
the existing transfer gives the requested asymptotic. This explains exactly
what each ladder below must finish.

\section{Route A: spend the slack on more admissible carriers}

\subsection{The relaxed graph and its cuts}

Use $\tau=1+\delta/2$, so $\gamma=\delta/2$. Join a source $a$ to a target $b$
when
\begin{equation}
 X\bigl(a/\gcd(a,b)\bigr)\le\tau.                        \label{edge}
\end{equation}
An edge supplies
$H_X(b)\ge X(a)-(1+\tau)$. Therefore a matching with deficiency $t$ satisfies
\eqref{finish} with $E=0$. The threshold one in the older graph was stronger
than the asymptotic goal requires.

\begin{lemma}[Upward-closed Hall cuts]
If $a\mid a'$, then $N(a')\subseteq N(a)$ in \eqref{edge}. A maximum-deficit
set $S\subseteq A$ is upward closed under divisibility in $[1,M]$.
Let $C=[1,M]\setminus S$ and $B=I\setminus N(S)$. Then $C$ is a divisor
down-set and
\begin{equation}
 |S|-|N(S)|=|B|-|C|.                                    \label{cut}
\end{equation}
For $b\in B$, $d\mid b$, $du\le M$, one has
\begin{equation}
 X(u)\le\tau\quad\Longrightarrow\quad du\in C.           \label{closure}
\end{equation}
\end{lemma}
\begin{proof}
Prime by prime, $a/\gcd(a,b)$ divides $a'/\gcd(a',b)$. This proves the
neighborhood inclusion. Adding an omitted multiple of an element of $S$
adds no neighbor and strictly increases the deficit. The full tail itself
is upward closed. The cardinality identity is immediate since both ambient
sets have size $M$. If $du\notin C$, then $du\in S$ and the divisor $d$
would supply an edge to $b$, proving \eqref{closure}.
\end{proof}

The open arithmetic assertion is now precise: prove
$|B|\le|C|+\varepsilon_\delta n$ for these cuts, retaining the original exact
rows \eqref{rows}, where $\varepsilon_\delta$ satisfies \eqref{finish}.
The stronger estimate $t=O_\delta(r)$, with $r=O_\delta(n/\log n)$ active rows,
would suffice, but is not necessary and is not proved.

\subsection{A bank can pay a whole cluster}

\begin{lemma}[One-bank margin]
For any $g$ with $X(g)>1$ and any set $S$ of distinct multiples of $g$ in
$[1,M]$,
\begin{equation}
 \sum_{\substack{b\in I\\g\mid b}}H_X(b)
 \ge\sum_{a\in S}(X(a)-2)+|S|.                          \label{cluster}
\end{equation}
\end{lemma}
\begin{proof}
Put $m=\lfloor M/g\rfloor$. The target quotients form an interval of $N\ge m$
consecutive integers; the source quotients are a distinct subset of $[1,m]$.
Every prime-power divisibility count on the target is at least the prefix
count. Hence the sum of the raw quotient loads on the target is at least
the corresponding source sum. Writing $s=X(g)>1$, the difference between
the left side and $\sum_S(X(a)-2)$ is at least
$|S|+(N-|S|)(s-1)\ge|S|$.
\end{proof}

Repeated cofactors are therefore not intrinsically a loss. The unsolved
part is allocating overlapping banks of different clusters. Pairwise cold
intersections are insufficient: with weights $0.4$ on $2,3,5,7$, the roots
$30,42,70,105$ have excess $0.2$ each, but their common target $210$ has
hinge only $0.6$. The exact multiway excess over the lcm is
\[
 \sum_i(X(d_i)-1)-(X(\operatorname{lcm}_i d_i)-1)
 =\sum_p w_p\left(\sum_i v_p(d_i)-\max_i v_p(d_i)\right)-(s-1).
\]
A proof must control this expression jointly, not merely its pairwise parts.

\subsection{A universal two-carrier reduction}

Let $K=\#\{b\in I:X(b)>1\}$. Every hot integer $u\le M$ satisfies
$\lfloor M/u\rfloor\le K$, since every target multiple of $u$ is hot.
Fix $\eta\le\delta/8$ and let $V$ contain all prices at most $\eta$.

For each source, if $V(a)\le1+\delta/4$, remove all its light and inactive
factors. The remaining carrier contains only prices $>\eta$, and its
complement has load at most $\tau$. If the complement is one, remove one
active prime as well; this keeps the carrier proper and admissible.
After the common load cutoff $X(a)\le L$, this carrier has at most $L/\eta$
active occurrences.

If $V(a)>1+\delta/4$, order the light occurrences by \emph{decreasing} prime
size and stop when their load first exceeds $t_0=1+\delta/4$. Their product
$u$ obeys
\[
 t_0<X(u)\le t_0+\eta\le1+3\delta/8<\tau.
\]
Remove this block and all inactive factors. The remaining carrier $h$ has
\begin{equation}
 h\le a/u\le\lfloor M/u\rfloor\le K,
 \qquad X(h)>1+5\delta/8.                               \label{dichotomy}
\end{equation}
If $p$ is the boundary light prime, the removed block has more than
$t_0/\eta$ occurrences, all at least $p$. Thus
\begin{equation}
 p<M^{\eta/t_0};
 \quad\text{every light prime left in }h\text{ is at most }p.
 \label{reverse}
\end{equation}
This is an unrestricted carrier reduction: either a pure substantial-price
carrier, or a numerically small carrier with a controlled residual light
support. It does not assert that either family already has a matching.
If $\eta$ is allowed to shrink with $M$, the degree bound $L/\eta$ grows;
one cannot subsequently invoke a fixed-degree certificate without a new
uniform estimate.

\subsection{An actual small core and a demand bound, with their limits}
In the sparse branch let $T=M/(K+1)\to\infty$. For a fixed sieve parameter
$U\ge8$, choose
\[
 \eta\le\min\{\delta/8,\ t_0\log T/(U\log M)\},\qquad z=T^{1/U}.
\]
Then every light prime remaining in a small carrier lies below $z$, and its
bank length is greater than $T=z^U$. This creates the length gap required by
the conditional sieve. The pure-heavy alternative has no diffuse root mass.

There is also a useful aggregate bound. Every integer at most $T$ is cold,
so for $p\le z$,
\[
 w_p\le\frac1{\lfloor\log T/\log p\rfloor}
 \le\frac U{U-1}\frac{\log p}{\log T}.
\]
Chebyshev's estimate and partial summation consequently give
\begin{equation}
 \sum_{p\le z}\frac{w_p}{p-1}\ll\frac1{U-1},\qquad
 \sum_{h\le K}V_{\le z}(h)\ll\frac K{U-1}.                \label{coremean}
\end{equation}
This counts distinct roots once. It is not an estimate for their source
multiplicities.

For sufficiently small $\eta$, a positive optimum also has
$\mu_V<1+\delta/8$: otherwise the earlier diffuse high-mean argument already
implies $J_2(X)\le0$. A high-light source assigned to $h$ has a distinct
quotient $u=a/h$ with $V(u)>1+\delta/4$. The raw prefix variance therefore
proves the labelled fibre estimate
\begin{equation}
 \frac{m_h^{\rm high}}{\lfloor M/h\rfloor}
 \le\frac{512\eta(1+\delta/8)}{\delta^2}=:q_\eta,
 \qquad \beta_h^{\rm high}\le q_\eta(X(h)-1).             \label{dilution}
\end{equation}
All inactive factors of $u$ are counted. If a root serves both branches,
\eqref{dilution} applies only to its high-light contribution. Summing these
small fractions independently at one target is still unjustified; joint
overlap is the open step.

\section{Route B: diffuse fluctuations in arbitrary arithmetic banks}

The analytic input is the fundamental lemma of the sieve for arbitrary
subsets of primes \cite{GS}. In a form used in the preceding checkpoint,
for primes $P\subseteq[2,z]$ and a real interval $J$ of length $N\ge z^u$,
\begin{equation}
 \#\{m\in J:(m,\prod_{p\in P}p)=1\}
 =(1+O(E(u)))N\prod_{p\in P}(1-1/p),\quad E(u)\to0,
 \label{sieve}
\end{equation}
uniformly in the shift and the subset. The usual additive sieve remainder
is absorbed into the relative error using $N\ge z^u$.

\begin{theorem}[Uniform one-sided bank concentration]
Let $V(t)=\sum_p\theta_pv_p(t)$, $0\le\theta_p\le\eta$, with arbitrary finite
support. Set $\mu_N=\sum_{p\le N}\theta_p/(p-1)\le B$. For every real $N\ge1$
and every real interval $J$ of length $N$,
\begin{align}
 \sum_{a\le N}|V(a)-\mu_N|&\le\sqrt{8\eta B}\,N,\label{sourcecon}\\
 \sum_{b\in J}(\mu_N-V(b))_+&\le\psi_B(\eta)N,\quad
 \psi_B(\eta)\longrightarrow0.                          \label{targetcon}
\end{align}
For sufficiently small $\eta$ depending only on $B$, one may take
\begin{equation}
 \psi_B(\eta)=2\eta^{1/4}
 +2B\exp\left(-\tfrac12\eta^{-1/4}+\tfrac12B\right).
 \label{psi}
\end{equation}
The support size, its reciprocal mass, the interval shift, and the number
of banks are unrestricted.
\end{theorem}
\begin{proof}
The raw prime-power moment calculation gives
$\sum_{a\le N}(V(a)-\mu_N)^2\le8\eta\mu_N\lfloor N\rfloor$; Cauchy--Schwarz
proves \eqref{sourcecon}. Primes exceeding $N$ vanish on the prefix and can
only improve the lower-deviation estimate on the target.

For the latter, put $u=\lceil\eta^{-1}\rceil$ and $z=N^{1/u}$. For $z\ge2$,
independently retain each supported prime $p\le z$ with probability
$1-e^{-t\theta_p}$. Averaging the upper bound in \eqref{sieve} over these
random subsets yields
\[
 \sum_{m\in J}\exp\left(-t\sum_{p\le z}\theta_p\ind_{p\mid m}\right)
 \le(1+E(u))N\prod_{p\le z}
       \left(1-\frac{1-e^{-t\theta_p}}p\right).
\]
This error is relative before averaging. Put $\mu_0=\sum_{p\le z}\theta_p/p$.
Chebyshev's prime bound and partial summation give
$0\le\mu_N-\mu_0\le d_\eta=O(\eta(1+\log u))$. Since
$1-e^{-x}\ge x-x^2/2$, the displayed product is at most
$\exp(-t\mu_0+t^2\eta B/2)$. Therefore
\[
 \frac1N\#\{m\in J:V(m)<\mu_N-\varepsilon\}
 \le(1+E(u))\exp(-t\varepsilon+td_\eta+t^2\eta B/2).
\]
Take $t=\eta^{-1/2}$, $\varepsilon=\eta^{1/4}$. For small $\eta$,
$d_\eta\le\varepsilon/2$ and $E(u)\le1$. The lower deficit is at most
$\varepsilon$ off this exceptional set and at most $B$ on it. There are
at most $2N$ integer positions, proving \eqref{psi}. If $z<2$, then
$\mu_N=O(\eta(1+\log u))$ and the estimate follows directly.
\end{proof}

No coprimality of a root and its quotient is needed: raw additivity gives
$V(dt)=V(d)+V(t)$. The correct baseline in the bank of $d$ is consequently
$V(d)+\mu_{M/d}$, not a common scalar independent of the root.

\subsection{A new unrestricted reduction at positive density}

At the genuine optimum, take any block $V$ of prices at most $\eta$, and
write $D_V=\sum_A V$, $\rho=n/M$. Exact active rows give
$D_V=\sum_Iz_bV(b)$. Equations \eqref{sourcecon}--\eqref{targetcon} imply
\[
 D_V\le\mu n+\sqrt{8\eta B_\delta nM},\qquad
 D_V\ge\mu F-\psi_{B_\delta}(\eta)M.
\]
Since $F\ge cn$, elimination of $\mu$ gives
\begin{equation}
 \boxed{\quad
 \frac{D_V}{n}\le
 \frac{c}{c-1}\sqrt{\frac{8\eta B_\delta}{\rho}}
 +\frac{\psi_{B_\delta}(\eta)}{(c-1)\rho}.
 \quad}                                                   \label{dense}
\end{equation}
Thus the small-price source demand vanishes uniformly as $\eta\to0$ on every
branch $\rho\ge\rho_0>0$, without a restriction on the prime profile.
The explicit factors of $\rho^{-1}$ show why this does not solve the sparse
branch by itself.

\subsection{Conditional cells without a bound on their number}

If the diffuse primes $P$ and excluded heavy primes $Q$ are disjoint subsets
of $[2,z]$, the same randomization retains the factor
$\mathfrak v_Q=\prod_{q\in Q}(1-1/q)$. On any quotient interval
$N\ge z^U$, it gives relative lower-deviation error
$o_{\eta,U}(N\mathfrak v_Q)$, uniformly in $Q$. The source absolute-error
estimate has the same scale. To retain every raw prime power, split at
$p^j=\sqrt N$; the high levels have relative contribution at most
\[
 2\eta z^2/\sqrt N\le2\eta\,2^{2-U/2}.
\]
The full proof is in \texttt{route\_core/CONDITIONAL\_DIFFUSE\_CELLS.md}.
This removes a cardinality restriction on the conditioned core; it requires
the genuine length gap $N\ge z^U$. It does not authorize conditioning on
arbitrary short cells.

\section{Joining the routes: one normalized and clipped target budget}

Retain $\tau=1+\delta/2$ and $\gamma=\delta/2$. For each source allow any
proper divisor $d\mid a$ with $X(a/d)\le\tau$. Choose fractions
$\theta_{a,d}\ge0$, $\sum_d\theta_{a,d}=1$, and set
\begin{align*}
 s_d&=\sum_a\theta_{a,d}\bigl(X(a)-(1+\tau)\bigr),&
 N_d&=\lfloor M/d\rfloor,&\beta_d&=s_d/N_d,\\
 S_\tau&=\sum_A(X(a)-(1+\tau)),&
 G(b)&=\sum_{d\mid b}\beta_d,&
 \mathcal N&=\sum_d\beta_dM/d.
\end{align*}
All actual sources are counted, including different sources with identical
active parts. Every root has $X(d)>1+\gamma$, and $N_d\ge2$. Hence
\begin{equation}
 \sum_I G\ge S_\tau,
 \qquad \mathcal N\le\tfrac32S_\tau=O_\delta(n).          \label{norm}
\end{equation}
The proof is elementary: each source is an actual multiple of its root;
each interval has at least $N_d$ such multiples; and
$(M/d)/\lfloor M/d\rfloor\le3/2$ when $N_d\ge2$.
There is no prime-count estimate for a collapsed fibre and no bound on
inactive cofactor multiplicity.

Let $V$ be any original price block with prices at most $\eta$, $W=X-V$,
and define
\[
 \nu_d=V(d)+\sum_{\substack{p\le M/d\\p\in\operatorname{supp}V}}
                  \frac{w_p}{p-1},\qquad
 \bar\nu(b)=\frac{\sum_{d\mid b}\beta_d\nu_d}{G(b)}\quad(G(b)>0).
\]
The model target hinge is $\widetilde H(b)=\pos{W(b)+\bar\nu(b)-1}$.
This auxiliary load is generally nonadditive. Original row balances and
optimality are \emph{not} transferred to it.

For a fixed $\Lambda\ge1$, use the accounting functional
\begin{equation}
 \mathcal C_\Lambda(H)=
 \sum_{G(b)>0}\min\{\Lambda G(b),H(b)\}
 +\sum_{G(b)=0}H_X(b).                                  \label{account}
\end{equation}
Increasing $\Lambda$ raises an accounting cap, not the target capacity;
each target still pays at most its one actual hinge.

\begin{theorem}[Joint demand-scale transfer]
Uniformly in all fractional choices, weights, shifts, and densities,
\begin{align}
 \mathcal C_\Lambda(H_X)&\ge\mathcal C_\Lambda(\widetilde H)
       -\frac{\Lambda\psi_{B_\delta}(\eta)}\gamma\mathcal N,
       \label{joint}\\
 J_c(X)&\le S_\tau-\gamma n
       -\mathcal C_\Lambda(\widetilde H)
       +\frac{\Lambda\psi_{B_\delta}(\eta)}\gamma\mathcal N.
       \label{mainnew}
\end{align}
In particular the error is $O_\delta(\Lambda\psi_{B_\delta}(\eta)n)$ even
when $n/M\to0$.
\end{theorem}
\begin{proof}
On the union of banks, $H_X(b)>\gamma$. A positive increase of the clipped
model payment can occur only where $\Lambda G>H_X$, hence $G>\gamma/\Lambda$.
The hinge is 1-Lipschitz and the positive part is convex, so that increase
is at most
\[
 \frac\Lambda\gamma\sum_{d\mid b}\beta_d(\nu_d-V(b))_+.
\]
For a multiple $b$ of $d$, raw additivity identifies the last deficit as
$\pos{\mu_{M/d}-V(b/d)}$. Sum \eqref{targetcon} in each bank, and then use
\eqref{norm}. This proves \eqref{joint} without separate capacity budgets.
Finally $J_c=S_\tau-\gamma n-\sum_IH_X$ and
$\mathcal C_\Lambda(H_X)\le\sum_IH_X$, proving \eqref{mainnew}.
\end{proof}

\paragraph{The exact unproved endpoint.}
It would suffice to find, for each $\delta$, a fixed $\Lambda_\delta$ and
admissible fractional assignments such that, uniformly for all sufficiently
small $\eta$ on the remaining genuine optimum class,
\begin{equation}
 \boxed{\qquad
 \mathcal C_{\Lambda_\delta}(\widetilde H)
 \ge S_\tau-\frac{\gamma n}{2}.
 \qquad}                                                  \label{open}
\end{equation}
Choose $\eta$ afterwards so the proved error is at most $\gamma n/4$;
then \eqref{mainnew} gives $J_c\le-\gamma n/4$.
An uncontrolled dependence $\Lambda=\Lambda(\eta)$ does not justify this
order. A quantified alternative would be enough if
$\Lambda(\eta)\psi_{B_\delta}(\eta)\to0$.
\textbf{Inequality \eqref{open} is not proved.} In particular, the retained
root values $V(d)$ cannot be erased or replaced by a common scalar.

\subsection{A stronger normalization that retains the cluster reserve}
A second construction uses the full credit in \eqref{cluster}. Choose the
first increasing active prime prefix $g(a)$ whose load exceeds
$1+\delta/2$. Then
\[
 1+\delta/2<X(g(a))\le2+\delta/2<X(a).
\]
An active prime remains after the root. Thus every active prime of $a$ is
at most its bank length $T_a=M/g(a)$. For a cluster $A_g$, set
\[
 m_g=|A_g|,\quad s_g=\sum_{A_g}(X(a)-2),\quad
 B_g=\sum_{b\in I:g\mid b}H_X(b),\quad\alpha_g=s_g/B_g.
\]
By \eqref{cluster}, $B_g\ge s_g+m_g$, so $\alpha_g<1$. With
$F_0(b)=\sum_{g\mid b}\alpha_g$ and $S_2=\sum_A(X(a)-2)$,
\begin{equation}
 \sum_IH_X(b)F_0(b)=S_2,\qquad
 \sum_g\alpha_gM/g\le\frac3\delta S_2=O_\delta(n).        \label{clusternorm}
\end{equation}
The coefficients retain the unit reserve per source. If $X(a)\le L$,
then $\alpha_g\le(L-2)/(L-1)$.

For a fixed $N_0$, roots with $M/g\le N_0$ are $N_0$-smooth. Their number is
at most $(1+\log_2M)^{\pi(N_0)}$, each has at most $N_0$ sources, and their
weighted source cost is
\[
 O_{N_0}((\log M)^{\pi(N_0)+1})=o(n),
\]
since $M<16n^2$. A diagonal cutoff makes retained bank lengths tend to infinity.
Hot-gcd mergers only increase lengths; their norm bound persists if the
new root load stays a fixed amount above one. Source deletions
must remain explicit in every subsequent balance identity.

Using the $\alpha_g$-weighted version of $\bar\nu$, define a new model hinge
$H_{\rm mod}$. The functional
\[
 \mathcal Q_\Lambda(H)=
 \sum_{F_0>0}\min\{\Lambda F_0,1\}H
 +\sum_{F_0=0}H_X
\]
again uses one target budget. The pointwise error is bounded by
$\Lambda\sum_{g\mid b}\alpha_g(\nu_g-V(b))_+$. Therefore
\begin{equation}
 J_c(X)\le S_2-\delta n-\mathcal Q_\Lambda(H_{\rm mod})
       +O_\delta(\Lambda\psi_{B_\delta}(\eta)n).          \label{clustertransfer}
\end{equation}
Proving $\mathcal Q_\Lambda(H_{\rm mod})\ge S_2-\delta n/2$ with the same
uniform parameter order is an alternative sufficient endpoint. It remains
unproved. This construction already pays a one-prime source family as one
cluster and avoids the artificial congestion diagnosed below.

There is a concrete all-occurrence subtarget. Put
$\mathcal V=\sum_IH_XF_0^2$. Since $\min(2t,1)\ge2t-t^2$,
$\mathcal Q_2(H_X)\ge2S_2-\mathcal V$. Its diagonal satisfies
\[
 \sum_g\alpha_g^2B_g
 \le S_2-\frac{\delta}{1+\delta}n;
\]
the off-diagonal is exactly
$2\sum_{g<h}\alpha_g\alpha_h B_{\operatorname{lcm}(g,h)}$.
This identifies a joint weighted lcm energy to attack. It is not bounded
by the diagonal reserve automatically, and an untruncated second moment
may overcharge rare targets with many incident roots.

For the full source set, the original rows also give the exact scale identity
\[
 \sum_{b\in I}z_bX_{>t}(b)
 =\sum_{a\in A:T_a>t}X_{>t}(a).
\]
Here $X_{>t}$ retains prime values $p>t$, including all their powers. If some
sources have been deleted, their demand is an additional term; it cannot
be erased from this identity. Controlling repeated use of this supply
across scales is a remaining stopping-certificate problem.

\section{Two checked arithmetic obstructions and their exact scope}

\subsection{Uniform routing can waste capacity even with one prime}

Take $w_2=1$, all other weights zero, $0<\delta<1$, and $I=[1,M]$.
The full tail is $A=\{v_2\ge3\}$, $n\sim M/8$. With active-part prime deletion
and payments $X(a)-2$, the root $2^j$ has limiting coefficient $(j-1)/4$.
At a target of valuation $t$,
\[
 G(b)\longrightarrow t(t-1)/8,\qquad H_X(b)=(t-1)_+.
\]
Therefore the unamplified congestion $E_X=\sum_I(G-H_X)_+$ obeys
\begin{equation}
 \frac{E_X}{M}\longrightarrow
 \sum_{t\ge9}\frac{(t-1)(t-8)}{8\,2^{t+1}}=\frac5{1024},
 \qquad \frac{E_X}{n}\longrightarrow\frac5{128}.           \label{oneprime}
\end{equation}
This refutes the blanket claim $E_X\le\delta n$ on arbitrary full tails
for $\delta<5/128$. It is not a positive optimum:
$J_c/M\to-1/4-\delta/8$. The unused target capacity matters.
For this limiting example, $H_X/G=8/t\le4$ wherever $H_X>0$; amplification
$\Lambda=4$ captures the whole hinge. A common first-crossing cluster also
pays the sources by \eqref{cluster}. The example diagnoses the routing
method, not the conjecture.

\subsection{Few bounded-degree hot roots can lose a linear union}

Let $M=2N$, $c=21/10$, $w_2=3/5$, and $w_p=1$ for
$\sqrt N<p\le N/2$, with all other prices zero. Take the actual full tail.
It contains all multiples of $16$, so $n\ge M/16-O(1)$, and every external
prime is used by the source $4p$. The roots $2p$ are minimally hot, have
load $8/5$ and degree two, and number only $O(N/\log N)=o(n)$.

Their prefix union, after division by two, consists of disjoint prime banks
on $[1,N]$, and has size $(\log2+o(1))N$. Average the quotient interval
$[C+1,C+N]$ over $C$ modulo the product of these primes. Its mean union size is
\[
 N\left[1-\prod_{\sqrt N<p\le N/2}(1-1/p)\right]
 =(1/2+o(1))N.
\]
Thus an actual shifted interval has union loss at least
$(\log2-1/2+o(1))N=\Omega(n)$. The usual conditions $K=O(n)$,
$r=O(n/\log n)$, prime sizes $O(n)$, bounded mean, and source exponential
tails all hold. They do not justify replacing an exponential union error
by the number of roots.

The exact active row, however, excludes this example. The source tail
consists of $v_2\ge4$, together with $v_2=2,3$ and one external occurrence.
Consequently
\[
 D_2(A)=\left(\frac58+\frac78\log2+o(1)\right)N
 =(1.2315\ldots+o(1))N.
\]
Every target multiple of four is hot, forcing 2-row supply at least
$(3/2+o(1))N$ in every shift. This contradicts \eqref{rows} for the positive
price $w_2$. The example is therefore a counterexample to the weaker
union-transfer inference, not to a lemma that truly retains optimality.

\section{The three detailed closing ladders and the order of attack}

\subsection*{A. Adaptive cuts, with cluster merging as the arithmetic tool}
\begin{enumerate}
\item Keep the true optimum and full tail. Spend half the fixed slack on
$\tau=1+\delta/2$, retaining $\gamma=\delta/2$ for the contradiction.
\item Use all admissible divisors, and reduce a worst Hall cut to the
upward-closed form \eqref{cut}--\eqref{closure}. The required error is the
fixed fraction in \eqref{finish}, not exact injectivity.
\item Use \eqref{cluster} to pay each common-hot-root cluster with its
unit-per-source margin. If cardinal matching is too restrictive, use
source demands $X(a)-(1+\tau)$ and actual capacities $H_X(b)$ instead.
\item Apply the universal carrier dichotomy \eqref{dichotomy}. Its low-light
branch has pure substantial-price carriers; its other branch has carriers
at most $K$. Both retain the original source multiplicities.
\item The genuinely open operation is to combine these banks while
controlling multiway overlap using \eqref{rows}. The bounded-degree root
counterexample proves that root count and $K=O(n)$ alone cannot perform it.
\item As soon as the adaptive deficit is below the fixed tolerance, apply
\eqref{finish}. There is no further integer-rounding obstacle.
\end{enumerate}

\subsection*{B. Joint bank transfer, followed by the model capacity inequality}
\begin{enumerate}
\item Normalize actual source demands as in \eqref{norm}. This already
resolves the uncontrolled sum of bank lengths and retains inactive factors.
\item Use the uniform one-sided lemma inside each bank, then clip only
after adding all contributions. This proves \eqref{mainnew} at scale $n$,
including the sparse branch. This step has been carried out.
\item Keep every root offset $V(d)$ and the unused target capacity in
\eqref{account}. The one-prime example rules out reverting to a stronger
unamplified congestion claim without an optimum-specific proof.
\item Prove \eqref{open} for a fixed amplification or a quantified rate
with $\Lambda\psi\to0$. The original exact prime rows remain available;
the auxiliary model has no new row equalities of its own.
\item Choose the diffuse cutoff only after the capture constants are
controlled. Then \eqref{mainnew} gives the strict contradiction immediately.
\end{enumerate}

\subsection*{C. An all-occurrence stopping certificate for the missing overlap}
\begin{enumerate}
\item Use the full-tail exponential moment to choose a finite source-load
cutoff $L$ at a prescribed weighted cost. Choose heavy and light price
cutoffs with the loss of any intervening band charged explicitly.
\item Form first-crossing hot clusters or the adaptive normalized banks.
Preserve their $O_\delta(n)$ weighted norm through every refinement. A
vanishing local relative error can then be summed at the correct scale.
\item If using heavy-profile cells, apply the conditional sieve only where
the actual quotient length is at least the necessary power of the largest
conditioned prime. Equations \eqref{reverse} can create a residual small
light core, but a shrinking price cutoff also increases the heavy depth.
\item Combine different arities within one signed or nonnegative target
certificate. A correction with a large lcm cannot be discarded just
because its prefix count is zero: it may occur at an actual target.
Track the signed norm and all such boundary terms jointly.
\item The remaining task is an arithmetic estimate for those simultaneous
corrections or, equivalently for this route, the capture assertion
\eqref{open}. Neither finite heavy degree nor a gap in active prices
removes inactive source cofactors. The invalid prime-count fibre argument
must not be reintroduced.
\end{enumerate}

The best immediate target is the relaxed adaptive cut or the weaker
amplified capture inequality, with exact original row balances. Further
separate three-, four-, or five-prime theorems would not discharge either
target. The positive-density branch now has the extra reduction
\eqref{dense}; the sparse branch has the genuinely demand-scale transfer
\eqref{mainnew}. The open arithmetic overlap estimate is common to both.

\section{Verification and limits of the evidence}

The displayed proofs were checked against raw prime powers, arbitrary shifts,
real quotient lengths, repeated sources within a carrier fibre, and the
distinction between actual and auxiliary optimality. The counterexample
calculations have independent derivations; exact rational checks of finite
instances are included. Small CRT-realizable searches tested intermediate
Hall statements and found no deficiency in the tested examples; these are
falsification probes, not proofs and not computations of new $g(n)$ values.

The report proves the new transfer and structural lemmas. It does not prove
the terminal Hall bound or \eqref{open}. The exact $2n$ conjecture is kept
separate: every closing implication here deliberately spends positive
$\delta$ slack.

\begin{thebibliography}{9}
\bibitem{WangRaw} W. Wang, raw fractional-to-integer transfer manuscript,
\url{https://github.com/science-prince/erdos-708-interval-divisor-covers/blob/main/papers/main-paper/raw-transfer.tex}.
\bibitem{WangMain} W. Wang, main interval-divisor cover manuscript,
\url{https://github.com/science-prince/erdos-708-interval-divisor-covers/blob/main/papers/main-paper/main.tex}.
\bibitem{GS} A. Granville and K. Soundararajan,
\emph{Multiplicative number theory: The pretentious approach},
Theorem 1.4.2 (fundamental lemma of sieve theory),
\url{https://dms.umontreal.ca/~andrew/PDF/Book.To1.4.pdf}.
\end{thebibliography}
\end{document}
